\section*{Ethics Statement}

This paper argues that a class of reported debiasing results overstates what
was achieved. We want to be careful about what follows from that and what does
not.

It does not follow that the interventions we decompose are harmful or that
their authors did anything improper. They report accuracy and bias score
exactly as the benchmark specifies, and BBQ's scaling of the ambiguous score
by the commitment rate is a deliberate design choice its authors motivate
explicitly. Our claim is about what the resulting number can support, not
about anyone's conduct. Nor does it follow that abstention is bad: declining
to answer an under-specified question about a person's group membership is
often the appropriate behavior, and a model that learns to do so has
learned something worth having. The problem we identify is narrower and
purely epistemic. The same movement that improves the score also widens the
identified set, weakening the evidence available to check whether the underlying disposition changed, so the
metric cannot distinguish a model that stopped preferring the stereotype from
one that merely stopped saying so.

That distinction has practical stakes, not only in principle. Qwen2.5-7B
posts a near-perfect first-party baseline score ($s_{\text{AMB}}=0.033$,
Table~\ref{tab:firstparty}) while its own stereotype-aligned rate $\theta$,
measured on the same items under forced choice, is 62.2\%. A deployment
decision made on the strength of the published-style score alone would be
relying on a number compatible, by the width of its own identified interval,
with a meaningfully stereotype-leaning disposition underneath. Our
recommendation is therefore additive rather than prohibitive: report the
abstention rate and the committed-answer bias alongside the headline score,
and report the identified interval, so readers can see how much the
evaluation actually pins down.

We work exclusively with existing public benchmarks (BBQ, WinoGender) and
public model checkpoints. No new human-subjects data was collected. These
benchmarks encode a US-centric, largely binary treatment of gender and a
specific inventory of social categories, and WinoGender carries documented
template defects \citep{gautam2024winopron}. Results here should not be read
as characterizing bias in any absolute sense, nor as transferring to other
languages, cultures, or deployment settings. Finally, our experiments
deliberately elicit stereotype-aligned outputs from models by removing their
ability to decline. This is done to measure disposition, and we release
per-item outputs for verification, not as a technique for eliciting harmful
content in deployment.
